\section*{Chapter \ref{ch:g8:powers} --- Powers}

\begin{solution}{exo:g8:powers:1}
$2^6 = 64$; \quad $3^3 = 27$; \quad $10^5 = 100\,000$; \quad
$1^{100} = 1$; \quad $0.1^2 = 0.01$; \quad $6^0 = 1$.
\end{solution}

\begin{solution}{exo:g8:powers:2}
$(-3)^2 = 9$; \quad $-3^2 = -9$; \quad $(-1)^{15} = -1$; \quad
$(-5)^3 = -125$; \quad $\left(\frac23\right)^2 = \frac49$.
\end{solution}

\begin{solution}{exo:g8:powers:3}
$7^4 \times 7^5 = 7^9$; \quad
$\dfrac{2^9}{2^3} = 2^6$; \quad
$\left(10^3\right)^4 = 10^{12}$; \quad
$5^6 \times 5^{-2} = 5^4$; \quad
$3^4 \times 7^4 = 21^4$ (same \omterm{def:g8:powers:def}{exponent}: multiply the bases).
\end{solution}

\begin{solution}{exo:g8:powers:4}
$10^{-2} = 0.01$; \quad
$4 \times 10^3 = 4000$; \quad
$2.5 \times 10^{-4} = 0.00025$; \quad
$10^0 = 1$.
\end{solution}

\begin{solution}{exo:g8:powers:5}
$10^5$; \quad $10^{-1}$; \quad $10^{10}$; \quad $10^{-6}$.
\end{solution}

\begin{solution}{exo:g8:powers:6}
$\dfrac{10^7 \times 10^{-3}}{10^2} = \dfrac{10^4}{10^2} = 10^2 =
100$.

$\dfrac{2^5 \times 2^4}{2^6} = \dfrac{2^9}{2^6} = 2^3 = 8$.

$\left(2^2\right)^3 \times 2^{-4} = 2^6 \times 2^{-4} = 2^2 = 4$.
\end{solution}

\begin{solution}{exo:g8:powers:7}
$2^3 \times 2^4 = 2^{12}$: \emph{false} --- \omterm{def:g8:powers:def}{exponents} add:
$2^7$.

$5^2 + 5^3 = 5^5$: \emph{false} --- no rule for \omterm{def:g1:addition:def}{sums}:
$25 + 125 = 150$, while $5^5 = 3125$.

$(3^2)^4 = 3^8$: \emph{true}.

$10^3 \times 10^3 = 100^3$: \emph{true} --- both equal $10^6$
(left: $10^{3+3}$; right: $(10^2)^3$).
\end{solution}

\begin{solution}{exo:g8:powers:8}
New people on day $4$: each of the $3^3 = 27$ people informed on day
3 tells three others: $3^4 = 81$. Knowing at the end of day 4:
$3 + 9 + 27 + 81 = 120$ people.
\end{solution}

\begin{solution}{exo:g8:powers:9}
\emph{1.} Thickness: $2^{10} \times 10^{-4}$ m
$= 1024 \times 10^{-4}$ m $\approx 0.1$ m $= 10$ cm.

\emph{2.} $2^{42} \times 10^{-4} \approx 4.4 \times 10^{12} \times
10^{-4} = 4.4 \times 10^8$ m, larger than $3.8 \times 10^8$ m: after
$42$ (theoretical!) folds, the wad of paper would pass the Moon.
\end{solution}

\begin{solution}{exo:g8:powers:10}
$2^{10} = 1024$; $10^3 = 1000$; $3^6 = 729$; $5^4 = 625$. Order:
\[
5^4 < 3^6 < 10^3 < 2^{10} .
\]
\end{solution}

\begin{solution}{exo:g8:powers:11}
$2^{100} = \left(2^{10}\right)^{10} = 1024^{10}$ and
$10^{30} = \left(10^3\right)^{10} = 1000^{10}$. Since $1024 > 1000$,
multiplying ten copies of each keeps the inequality:
$2^{100} > 10^{30}$.
\end{solution}

\begin{solution}{pb:g8:powers:1}
\textbf{1.} The grains double from square to square starting at
$1 = 2^0$: square $k$ holds $2^{k-1}$ grains
(\cref{def:g8:powers:def}). Square $64$ holds $2^{63}$.

\textbf{2.} $S_1 = 1$, $S_2 = 1 + 2 = 3$, $S_3 = 3 + 4 = 7$,
$S_4 = 7 + 8 = 15$, $S_5 = 15 + 16 = 31$: always one less than
the next power of $2$ ($2$, $4$, $8$, $16$, $32$). Conjecture:
$S_n = 2^n - 1$.

\textbf{3.} Doubling every term of $S_n$ shifts each power up by
one ($2 \times 2^{k} = 2^{k+1}$):
\begin{align*}
S_n &= 1 + 2 + 4 + \dots + 2^{n-1}, \\
2 \times S_n &= \phantom{1 + {}} 2 + 4 + \dots + 2^{n-1} + 2^n .
\end{align*}
Subtracting the first line from the second, every term from $2$
to $2^{n-1}$ appears in both and cancels:
\[
2 \times S_n - S_n = 2^n - 1,
\qquad\text{that is}\qquad
S_n = 2^n - 1 .
\]

\textbf{4.} The total is $S_{64} = 2^{64} - 1$ grains. A
sixty-fifth square would hold $2^{64}$ grains: the whole board
carries exactly one grain fewer than that single square would.

\textbf{5.} The first \omterm{def:g3:division:half}{half} holds $S_{32} = 2^{32} - 1$ grains.
The whole board holds $2^{64} - 1$, so the second \omterm{def:g3:division:half}{half} holds
\[
\left(2^{64} - 1\right) - \left(2^{32} - 1\right)
= 2^{64} - 2^{32}
= 2^{32} \times \left(2^{32} - 1\right)
\]
(factor out $2^{32}$, using $2^{32} \times 2^{32} = 2^{64}$,
\cref{thm:g8:powers:rules}): exactly $2^{32}$ times the first
\omterm{def:g3:division:half}{half}.

\textbf{6.} By the \omterm{def:g8:powers:def}{exponent} rules, $2^{64} = 2^{4 + 60} =
2^4 \times \left(2^{10}\right)^6$. Since $2^{10} = 1024 > 10^3$,
\[
2^{64} > 16 \times \left(10^3\right)^6 = 16 \times 10^{18}
= 1.6 \times 10^{19} .
\]

\textbf{7.} More than $1.6 \times 10^{19}$ grains at
$5 \times 10^{-2}$ g each:
\[
1.6 \times 10^{19} \times 5 \times 10^{-2}
= 8 \times 10^{17} \text{ g} .
\]
Dividing by $10^6$ g per tonne: more than $8 \times 10^{11}$
tonnes --- eight hundred billion tonnes.

\textbf{8.} $\dfrac{8 \times 10^{11}}{8 \times 10^{8}} = 10^3$:
the king promised at least a \emph{thousand years} of today's
entire world harvest. (The legend says his advisers told him as
much.)

\textbf{9.} $2^{19} = 2^9 \times 2^{10} = 512 \times 1024 =
524\,288 < 10^6$, while $2^{20} = \left(2^{10}\right)^2 =
1024^2 = 1\,048\,576 > 10^6$. So the smallest \omterm{def:g8:powers:def}{exponent} is
$n = 20$: twenty doublings pass the million.

\textbf{10.} The pay on day $n$ is $2^{n-1}$ cents (day 1:
$2^0 = 1$). Now
\[
2^{26} = 2^6 \times \left(2^{10}\right)^2
= 64 \times 1\,048\,576 = 67\,108\,864 < 10^8,
\]
\[
2^{27} = 2 \times 2^{26} = 134\,217\,728 > 10^8 .
\]
So the daily pay first exceeds $10^8$ cents when $n - 1 = 27$:
on day $28$.

\textbf{11.} $21 = 16 + 4 + 1$: weights $16$, $4$ and $1$ g.
$27 = 16 + 8 + 2 + 1$: weights $16$, $8$, $2$ and $1$ g.

\textbf{12.} By question 3, the weights $1, 2, 4, 8$ together
weigh $S_4 = 2^4 - 1 = 15$ g. So without the $16$ g weight the
merchant cannot pass $15$ g: any target of $16$ g or more must
use it. And a target of $15$ g or less must not use it, since the
$16$ g weight alone already exceeds the target. The choice of the
largest weight is therefore \emph{forced}. What remains is a
target of at most $15$ g to be formed with $1, 2, 4, 8$ --- and
the same argument repeats: $8$ is forced (used if the remaining
target is $\geq 8$, unused otherwise, because $1 + 2 + 4 = 7$),
then $4$ (because $1 + 2 = 3$), then $2$, then $1$.

\textbf{13.} Following the forced choices, the remaining target
after each stage is at most the total of the remaining weights,
so the process ends with \omterm{def:g3:division:remainder}{remainder} $0$: every target from $1$ to
$31$ is reached. And since every choice along the way was forced,
no other selection of weights can reach the same target: the
writing of each number from $1$ to $31$ as a \omterm{def:g1:addition:def}{sum} of distinct
powers of $2$ exists and is \emph{unique}.

\textbf{14.} The six weights total $S_6 = 2^6 - 1 = 63$ g, and
the same forced-choice argument covers every target from $1$ to
$63$ g. For $45$: the target is $\geq 32$, so use $32$; remains
$13 < 16$, so skip $16$; use $8$ (remains $5$), use $4$ (remains
$1$), skip $2$, use $1$:
\[
45 = 32 + 8 + 4 + 1 = 2^5 + 2^3 + 2^2 + 2^0 .
\]

\textbf{15.} Choosing a $0$ or a $1$ on each of the $64$ squares
amounts to choosing which powers $2^0, 2^1, \dots, 2^{63}$ to
include in a \omterm{def:g1:addition:def}{sum} --- exactly the merchant's weighing with $64$
weights. The smallest storable number is $0$ (all squares at
$0$); the largest is the \omterm{def:g1:addition:def}{sum} of \emph{all} the powers, which is
the king's total: $S_{64} = 2^{64} - 1$ (Part I). By the
forced-choice argument, every whole number in between is reached
exactly once: a $64$-bit machine stores precisely the whole
numbers from $0$ to $2^{64} - 1$.
\end{solution}
